\documentclass[12pt]{article}
\usepackage[a4paper,margin=1in]{geometry}
\usepackage{amsmath, amssymb, booktabs, setspace}
\usepackage{enumitem}
\usepackage{graphicx}

\title{Practice Problem Set 4\\ \vspace{10pt} \large ECON 101 -- Macroeconomics}
\author{}
\vspace{-10pt}
\date{ \vspace{-30pt} Fall 2025}

\begin{document}
\maketitle





\begin{enumerate}[label=\textbf{\arabic*.}, leftmargin=1.5em]

\item Explain the three main sources of long-run economic growth. Then, using Canada as an example, describe one policy that could enhance each source.

\vspace{1em}

\item Suppose the production function of a small open economy is:
\[
Y = A K^{0.3} L^{0.7}
\]
where $A$ is total factor productivity (TFP). If $A = 5$, $K = 400$, and $L = 100$:

\begin{enumerate}[label=(\alph*)]
\item Calculate the level of real GDP.
\item If capital increases by 10\% and labour remains constant, what is the percentage increase in output? (Use the exponent rule.)
\item Explain why output increases by less than 10\%.
\end{enumerate}

\vspace{1em}

\item A country’s real GDP grows by 4\%. Capital grows by 2\%, labour by 1\%, and the capital and labour shares are 0.4 and 0.6 respectively.

\begin{enumerate}[label=(\alph*)]
\item Compute the contribution of capital, labour, and total factor productivity (TFP) to growth.
\item Interpret what a negative TFP growth would imply for an economy.
\end{enumerate}

\vspace{1em}

\item Discuss how the following policies can affect long-run growth:
\begin{enumerate}[label=(\alph*)]
\item Increasing research and development (R\&D) subsidies
\item Expanding higher education
\item Strengthening property rights and reducing corruption
\end{enumerate}

\end{enumerate}

\newpage


\begin{enumerate}[label=\textbf{\arabic*.}, start=5, leftmargin=1.5em]

\item 
Write down the components of aggregate expenditure (AE). Which component is the largest in the Canadian economy? Explain how \textit{planned investment} differs from \textit{actual investment}.

\vspace{1em}

\item Consider the following data for a closed economy (all values in billions of dollars):

\begin{center}
\begin{tabular}{ll}
\toprule
\textbf{Component} & \textbf{Function / Value} \\
\midrule
$C = 100 + 0.8Y_D$ & where $Y_D = Y - T$ \\
$I = 150$ & \\
$G = 200$ & \\
$T = 100$ & \\
\bottomrule
\end{tabular}
\end{center}

\begin{enumerate}[label=(\alph*)]
\item Write the aggregate expenditure function $AE = C + I + G$.
\item Solve for equilibrium real GDP ($Y = AE$).
\item What is the marginal propensity to consume (MPC)?
\item Calculate the multiplier.
\item Suppose government spending increases by \$50 billion. Compute the new equilibrium GDP.
\end{enumerate}

\vspace{1em}

\item Explain what happens to unplanned inventories when:
\begin{enumerate}[label=(\alph*)]
\item Actual output $>$ planned AE
\item Actual output $<$ planned AE
\end{enumerate}
How do firms respond in each case to return the economy to equilibrium?

\vspace{1em}

\item Suppose net exports are given by:
\[
NX = 50 - 0.1Y
\]
Using the same parameters from Question 6, find the new equilibrium level of GDP. Explain intuitively how the \textit{marginal propensity to import} affects the multiplier.

\end{enumerate}


\end{document}
